%% kapa: exam=1 round=35 year=2024 period=1 form=A subject=economics subject_ko=경제학원론 seq=9 number=49
%% source: Q-Net 감정평가사 기출문제 (cst00309) artlSeq=5213229 file=raw/1008/5213229/1교시 문제지 원본.pdf
%% conversion: kordoc 4.24.51 -> tools/gen1.py (auto); review=proofread-sonnet; flags=none
\begin{question}{49}
두 재화 $X$, $Y$에 대해 양(+)의 소득 $M$을 가지고 효용을 극대화하는 갑의 효용함수는 $U(X,Y)=X+Y$이다. $Y$재 가격은 6이며, $X$재 가격은 5에서 8로 상승하였다. 이에 관한 설명으로 옳은 것은?
\begin{choices}
  \item $X$재 수요량 변화는 대체효과에 기인한다.
  \item $X$재 수요량 변화는 소득효과에 기인한다.
  \item $Y$재 수요량 변화는 없다.
  \item 수요량 변화의 1/3은 대체효과에 기인한다.
  \item 수요량 변화의 2/3는 소득효과에 기인한다.
\end{choices}
\end{question}
\answer{1}
